%!TEX root = ../thesis.tex
%*******************************************************************************
%****************************** Fourth Chapter **********************************
%*******************************************************************************
\part{Applications in Environmental Science}
\chapter{Shared Dynamics between Temporally Lagged Climate Variables: A Case Study of SST and Continental Precipitation}


% **************************** Define Graphics Path **************************
\ifpdf
    \graphicspath{{content/chapter4/figs/raster/}{content/chapter4/figs/pdf/}{content/chapter4/figs/}}
\else
    \graphicspath{{content/chapter4/figs/vector/}{content/chapter4/figs/}}
\fi

\section{Teleconnections in the Climate System}
complicated climate system, importance of teleconnections for understanding, predictions etc.; lagged teleconnections challening to analyze in a data-driven way, requires 
advanced techniques (complex rotated MCA)


\section{Case Study 1: Representation of the Madden-Julian Osciallation}
sub-seasonal time scales 

\section{Case Study 2: Shared Dynamics between SST and Precipitation}
seasonal to decadal time scales, large-scale patterns, regional impacts, etc.

Description of the datasets used: ERA5 dataset for sea surface temperature (SST) and continental precipitation; Time frame: 1960 to 2019; Geographic coverage: 40°S to 60°N, 1°x1° spatial resolution

Data Preprocessing Steps: temporal smoothing; normalizeation, weighting, etc.

Identification of large scale patterns: Seasonal Cycle; El Niño-Southern Oscillation (ENSO), Detection of ENSO-related patterns, spatial distribution of 
SST anomalies and their correlation with precipitation changes; Implications for predicting ENSO impacts on global weather patterns;
Global Warming Trend; Identification of trends in SST and precipitation linked to global warming; Regional analysis of warming effects and 
their correlation with precipitation changes; Pacific Decadal Oscillation (PDO); analysis of PDO patterns and their influence on climate variability;
Temporal lag and phase relationships between SST and precipitation during PDO events; ENSO Modoki; Detection of ENSO Modoki patterns; Differences from 
canonical ENSO and their specific impacts on regional climates; Atlantic Meridional Mode (AMM); Identification of AMM-related variability; Correlation between 
SST in the Atlantic and precipitation in adjacent continental regions; 

Key Findings and Insights:
Enhanced Understanding of Teleconnections; Insights into long-distance correlations between ocean and atmospheric variables; Identification of key regions and time lags 
in climate interactions; Improved Predictive Capabilities; Contributions to more accurate climate models by incorporating lagged relationships; Potential applications 
in seasonal and long-term climate forecasting; 

Regional Climate Insights
Detailed analysis of regional climate phenomena and their broader implications
Identification of regions with strong seasonal and interannual variability

Conclusion
Summary of key findings from the application of complex rotated MCA on climate variables
Implications for future climate research and potential areas for further study
Final thoughts on the importance of advanced data analysis techniques in enhancing our understanding of climate systems